\section{What the certificates support}

The theorem constructs simultaneous intervals for heterogeneous signed
moments when exact supported exposure laws, a sufficient conditional graph
and legal residual moments are supplied. Its practical resolution depends
on the particular query. A small active region can avoid irrelevant global
weights and graph connections, as the first-half E2 comparison demonstrates.
A small subgroup also contains less relevant evidence, and a small policy
step scales its signal along with its radius. Neither locality nor clipping
guarantees useful decisions in every supported design.

The completed E1 grid separates these issues empirically. The coverage
diagnostics pass throughout, including every heavy-tail cell with zero
detected violations. Stronger bounded inference recovers detection in
settings where Hoeffding or pure truncation does not. Consequently an
uninformative but valid interval should be interpreted as insufficient
evidence under its envelope and graph budget. It does not establish that
the query is safe, or that every valid method must be equally uninformative.
The scoped lower bound identifies an obstruction through subgroup mass,
exposure and clique replication; its expected-risk conclusion does not turn
our sufficient radii into exact information limits.

The same local information is available to the strongest matched comparator.
Its equality with COMET locates the contribution in the signed-moment
formulation, target-specific evidence-cost analysis, scoped limits and
decision contract. The experiments do not support a separate estimator
superiority claim over a fully informed generic graph method. Catoni and
OPE comparisons are restricted to their theorem-compatible submodels;
the differing guarantees and assumptions remain part of the comparison.

Independent discovery permits a random finite audit family to be frozen
before certification. It does not protect all unsearched subgroups.
Similarly, finite-policy certification controls return and unsafe under
the specified risk family. It does not promise frequent return, improved
utility, conditional risk at the same $\delta$, or future deployment safety.
The complete-pipeline results must therefore report baseline returns,
nonbaseline utility and service alongside the joint safe-and-useful event.
The E4 gap study further separates primary-query power from family-any
discovery. Its tested heavy-tail and weak-overlap environments provide no
primary-query certification. E5 remains in progress in this revision; its
partial outcomes do not enter the conclusions here.

The source graph and moment envelope are substantive model inputs.
Their declarations cannot rule out a hidden common shock or certify an
unknown tail scale from one evaluation log. The numerical implementation
also remains distinct from a formally verified outward-rounded realization
of the real-arithmetic proof. These boundaries define the experiments'
scope. The blocked ACS extension leaves real-covariate validation absent;
it does not supply evidence against or in favor of a causal social
intervention.

\section{Conclusion}

Counterfactual moment certification depends on the signed query, supported
residual moments and a sufficient graph of the randomness entering its
score. Explicit clipping bias and graph concentration yield a certificate
without equal row means or repeated independent macro-blocks. Query-local
analysis identifies the relevant costs, and direct contrasts preserve the
cancellation in a policy comparison. The empirical validity grid shows both
useful bounded inference and severe limits in the tested heavy-tail regimes.
Decision evidence must meet a further standard: the full risk and utility
family must support safe, useful return under fair strong comparisons.
